\chapter{Обзор методов передачи игровых данных}

\section{Надёжная и ненадёжная доставка данных}

В сетевых играх используются два основных подхода к передаче пакетов: надёжная доставка\cite{kurose2013networking} и доставка
без гарантий получения. Они различаются по задержке, стабильности и объёму служебной информации.

Надёжная доставка, основанная на протоколе TCP\cite{stevens2011tcpip}, обеспечивает контроль целостности данных и сохраняет 
порядок передачи пакетов, повторно отправляя утраченные сегменты. Но такая модель добавляет задержку из-за механизма подтверждений и повторных
передач, поэтому она не подходит для динамичных игр с высокой частотой обновления состояния.

Ненадёжная доставка, реализованная через протокол UDP, не гарантирует получение каждого отдельного пакета, но
обеспечивает минимально возможную задержку. Большинство сетевых игр используют UDP и свои собственные
механизмы, которые обеспечивают надёжность для критичных данных, что сохраняет эффективность
в условиях различных сетевых условий.

\section{Классификация и приоритизация игровых пакетов}

Игровые данные по-разному реагируют на задержки и потери, поэтому пакеты обычно разделяют на классы
приоритетности. Это позволяет оптимизировать передачу и избежать перегрузки канала связи.

Критичные данные включают команды игрока, результаты столкновений и важные события игровой логики.
Они должны доставляться максимально быстро и целостно, поскольку их потеря
может привести к заметной рассинхронизации состояния.

Важные, но допускающие потерю данные могут передаваться чаще, но без гарантии получения. Потеря части пакетов компенсируется
новыми обновлениями в следующих сообщениях.

Второстепенные данные (визуальные эффекты, косметические параметры) могут передаваться с
низким приоритетом или не передаваться вовсе, если клиент может их корректно воспроизводить локально.

\section{Механизмы обеспечения частичной надёжности}

Поскольку UDP сам по себе не предоставляет гарантированную доставку, игровые движки внедряют
дополнительные механизмы частичной надёжности, позволяющие гарантировать передачу важных
данных.

К использующимся приёмам относятся:

\begin{itemize}
    \item подтверждение доставки для отдельных типов пакетов;
    \item повторная отправка важных данных при отсутствии подтверждения;
    \item нумерация пакетов для восстановления порядка исполнения команд;
\end{itemize}

Эти механизмы позволяют экономно расходовать сетевой ресурс, сохраняя критичные данные в
актуальном состоянии.

\section{Адаптивные алгоритмы передачи данных}

Современные сетевые игры используют динамические методы, которые подстраиваются под качество канала.
Адаптивные алгоритмы позволяют снижать нагрузку в условиях высокой задержки или потери пакетов и
увеличивать точность передачи при стабильном соединении.

К распространенным техникам относятся:

\begin{itemize}
    \item изменение частоты отправки пакетов в зависимости от текущей задержки;
    \item снижение точности квантования при перегрузке сети;
    \item динамическое отключение второстепенных данных;
    \item регулирование частоты отправки обновлений в зависимости от важности объектов.
\end{itemize}

Благодаря таким методам происходит обеспечение более устойчивой работы сетевой игры при различных сетевых условиях, снижая риск
возникновения рассинхронизации.
